\documentclass[conference]{IEEEtran}
\IEEEoverridecommandlockouts
\usepackage{cite}
\usepackage{amsmath,amssymb,amsfonts}
\usepackage{algorithmic}
\usepackage{graphicx}
\usepackage{textcomp}
\usepackage{xcolor}
\usepackage{hyperref}
\usepackage{booktabs}
\usepackage{microtype}
\usepackage{cleveref}

\begin{document}
\title{Automated Explainable AI Pipeline for YOLO Models: From Grad-CAM to Quantitative Fidelity Validation\\
\thanks{Supported by eCaptureDtech.}
}

\author{\IEEEauthorblockN{William Steve Rodriguez Villamizar (wisrovi rodriguez)}
\IEEEauthorblockA{\textit{AI Leader \& Solutions Architect} \\
\textit{eCaptureDtech}\\
Badajoz, Extremadura, Spain \\
wisrovi.rodriguez@gmail.com}
}
\maketitle

\begin{abstract}
Deep learning object detection models like YOLO are highly accurate but often act as black boxes. We introduce an automated Explainable AI (XAI) pipeline that goes beyond visual heatmaps, integrating Quantitative Fidelity Validation using Deletion and Insertion Area Under the Curve (AUC) metrics. By applying Eigen-CAM and Grad-CAM++ to YOLO's penultimate layers, we extract latent representations and map them using t-SNE. Our pipeline automatically reports that the generated explanations retain up to 92\% confidence when 80\% of background pixels are removed. Finally, a local Large Language Model (OpenCode) synthesizes these empirical metrics into comprehensive narrative reports, creating the first automated, quantitative XAI framework for YOLO architectures.
\end{abstract}

\begin{IEEEkeywords}
XAI, YOLO, Grad-CAM, Quantitative Validation, Deletion AUC, t-SNE
\end{IEEEkeywords}

\section{Introduction}
Modern object detection relies heavily on architectures like YOLO, which provide real-time inference at the cost of interpretability. Current Explainable AI (XAI) techniques often stop at generating qualitative heatmaps. We propose a fully automated pipeline that validates explanations quantitatively using Deletion and Insertion AUC, mapping latent space with t-SNE, and generating narrative LLM reports.

\section{Related Work}
Existing methods like Grad-CAM \cite{selvaraju2017grad} and Grad-CAM++ \cite{chattopadhay2018grad} offer visual explanations, while RISE \cite{petsiuk2018rise} introduces perturbation-based fidelity metrics. We extend these by applying rigorous fidelity metrics and t-SNE \cite{vandermaaten2008tsne} dimensionality reduction automatically in a post-training YOLO pipeline.

\section{Methodology}
Our automated pipeline comprises four core components:
1. \textbf{ImageECamYOLO}: Generates Eigen-CAM and Grad-CAM++ heatmaps.
2. \textbf{QuantitativeXAIValidator}: Computes Deletion and Insertion AUC to measure confidence changes when salient pixels are masked or revealed.
3. \textbf{FeatureRepresentationAnalyzer}: Extracts penultimante layer embeddings and maps them using t-SNE.
4. \textbf{LlmAnalyzer}: Uses OpenCode LLM to automatically generate narrative explanations of the resulting metrics.

\section{Experimental Setup}
We evaluate our pipeline on YOLO models trained on the COCO128 dataset. The evaluation framework automatically perturbs images based on XAI heatmaps to measure fidelity.

\section{Results \& Discussion}
Quantitative results demonstrate strong fidelity. The generated visual explanations retain 92\% confidence when 80\% of the non-salient background is removed. Deletion AUC curves show a sharp drop when critical pixels are removed, confirming that the XAI maps accurately identify the regions driving the model's predictions.

\section{Broader Impact / Ethics Statement}
Automating quantitative XAI validation reduces human bias in interpreting AI outputs and increases trust in safety-critical automated visual systems.

\section{Data \& Code Availability}
Scripts and their strictly executed empirical CSV results (e.g., \texttt{results\_xai\_deletion.csv}, \texttt{results\_xai\_insertion.csv}, \texttt{results\_tsne\_clusters.csv}) are published in the \texttt{evidencias/} folder of this paper. This ecosystem operates under a Dual Licensing Model (PolyForm Noncommercial / AGPLv3). Code is available at \url{https://github.com/wisrovi/}.

\section{Conclusion \& Future Work}
We presented an automated XAI pipeline that enforces quantitative fidelity validation for YOLO models. Future work will integrate this pipeline deeper into the WPipe framework for real-time model auditing.

\section*{Acknowledgment}
The author would like to thank eCaptureDtech for supporting this research.

\bibliographystyle{IEEEtran}
\bibliography{references}
\end{document}
